\chapter{Final Conclusion and Future Improvements}

\section{Final conclusion}
\label{sec:final-conclusion}
\begin{sloppypar}
Chapter 1 opened this report with a fragmented Tunisian news landscape -- \textit{La Presse}, \textit{Business News}, \textit{Kapitalis}, and \textit{Mosaïque FM} publishing independently, with no shared index, no cross-source view, and no machine-readable signal attached to any of it -- and four concrete objectives to close that gap: O1, automated collection; O2, NLP enrichment; O3, GenAI access; and O4, reproducible MLOps (Table~\ref{tab:tech-stack}). Ten chapters later, the platform this report has traced end to end meets all four, each to a degree this closing chapter can now state precisely rather than as an aspiration. O1 runs today through \texttt{RSSScraper} polling four sources every time the pipeline executes, with \texttt{KapitalisScraper}'s deeper HTML crawl built, tested, and deliberately left dormant (Chapter 3, Section~\ref{sec:kapitalis}). O2 runs as a single batched pass over every new article -- multilingual e5 embeddings, XLM-RoBERTa sentiment, a ten-category LLM theme, and a gazetteer region tag -- with BERTopic clustering as a separate, operator-triggered stage over the same embeddings (Chapter 5). O3 is live in three places a reader actually touches: cached on-demand summaries, LLM-assisted search ranking, and a five-source retrieval-augmented chat assistant (Chapter 6). O4 is real but partial, and this report has been explicit about exactly where: a Prefect flow and a logged MLflow experiment exist and work, but the registered daily deployment has no worker to execute it, and a fifth, supporting objective -- a web dashboard usable by someone who does not read code -- is fully delivered as the six-route React frontend of Chapter 10. None of these four assessments is asserted from a plan or a design document; each is read directly from the code, the configuration, and the commit history this report has quoted throughout, including the places where what was actually built falls short of what was intended.
\end{sloppypar}

\section{Project contributions}
\label{sec:contributions}
\begin{sloppypar}
Layer by layer, the platform delivers a working, end-to-end pipeline rather than a set of disconnected proofs of concept. The \textbf{collection} layer (Chapter 3) contributes a shared \texttt{BaseScraper} template implemented twice -- an HTML listing crawler and an RSS collector -- converging on one validated \texttt{Article} shape, with layered error handling (per-request, per-article, per-source) and a durable, database-enforced uniqueness guarantee on \texttt{url}. The \textbf{storage} layer (Chapter 4) contributes a single PostgreSQL instance, extended with pgvector, holding both an article's relational metadata and its 384-dimension embedding in one table, indexed for the exact access patterns the API and dashboard exercise, and grown to its current shape through two small, additive migrations rather than a heavyweight migration framework. The \textbf{machine learning} layer (Chapter 5) contributes three model families working together -- multilingual sentence embeddings, fine-tuned sentiment classification, and unsupervised BERTopic clustering -- each wrapped by a lighter, explicitly-fallback-guarded heuristic (an LLM theme classifier, an LLM topic labeller, a gazetteer-first region tagger) so a single model failure degrades one column's quality rather than the whole run. The \textbf{GenAI} layer (Chapter 6) contributes a single fifty-line provider abstraction, \texttt{etl/llm.py}, underlying three reader-facing features -- cached article summaries, precision-tuned semantic search, and a five-source "ask the news" chat assistant -- all swappable to a different large language model provider through three environment variables. The \textbf{orchestration} layer (Chapter 7) contributes a scheduled, retried Prefect flow and a fully populated MLflow experiment for every clustering run, with a documented reliability posture that lets side channels -- a metrics row, a tracking log -- fail quietly while the steps that produce data fail loud. The \textbf{backend} layer (Chapter 8) contributes a five-router FastAPI surface, thirteen endpoints tabulated for the first time in this report, generating its own interactive documentation with no separate step. The \textbf{deployment} layer (Chapter 9) contributes a five-service Docker Compose stack, two of its five images built from this project's own source, brought up with a single command either as a fully containerized stack or as a lighter host-plus-database workflow. The \textbf{frontend} layer (Chapter 10) contributes six routed views sharing a persistent navigation header and a floating chat assistant, including a hand-built Tunisia governorate choropleth requiring no third-party map library at all. Across every one of these eight layers, this report has also named real, documented bugs found and fixed against a running system -- a scraper mis-selecting the wrong feed field, a search-ranking model padding irrelevant results, a database connection pinned idle across an LLM call, a flex-layout interaction breaking both scroll and images in the same panel -- evidence that News Insight was iterated against its own behaviour rather than written once from a specification and left alone.
\end{sloppypar}

\section{Future improvements}
\label{sec:future-improvements}
\begin{sloppypar}
Every chapter in this report has been as explicit about what is not yet done as about what is. Collected here, the concrete, verified gaps this platform still carries define the next iteration's backlog:
\end{sloppypar}

\begin{sloppypar}
\begin{itemize}
    \item \textbf{Broader source coverage.} Collection today is limited to the four RSS feeds already named -- \textit{lapresse}, \textit{businessnews}, \textit{kapitalis}, and \textit{mosaique} (\texttt{RSS\_FEEDS}, \texttt{etl/extract.py}) -- with no other Tunisian outlet wired in. \texttt{KapitalisScraper}'s deeper HTML crawl (Chapter 3, Section~\ref{sec:kapitalis}) is fully implemented and tested but sits behind a commented-out line in \texttt{build\_scrapers}; reactivating it, and writing equivalent listing crawlers for additional outlets, is the most direct way to widen the corpus beyond what RSS alone exposes.
    \item \textbf{A worker for the scheduled pipeline.} Chapter 7 (Section~\ref{sec:scheduling}) traced a real operational gap rather than a hypothetical one: \texttt{flows/schedule.py} registers a \texttt{daily-news-pipeline} deployment against a cron schedule, but no Prefect worker or work pool is defined anywhere in \texttt{docker-compose.yml} to actually execute it, so today the pipeline only runs through the two direct paths -- a manual CLI invocation or \texttt{POST /pipeline/trigger}. Adding a worker service to the Compose stack, polling the registered work pool, would let the deployment already defined in code run unattended for the first time.
    \item \textbf{An automated migration runner.} Chapter 9 (Section~\ref{sec:persistence}) showed that \texttt{db/schema.sql} only ever runs once, on an empty \texttt{pgdata} volume, through the official Postgres image's initialization convention; the two additive files under \texttt{db/migrations/} have no Compose service, Makefile target, or documented step that applies them against an already-running database, leaving that step to an undocumented manual \texttt{psql} session. A small migration-runner service, or a lightweight tool such as Flyway or a custom versioned-migrations script wired into the Makefile, would close that gap.
    \item \textbf{Uniform connection handling in the API.} Chapter 8 (Section~\ref{sec:data-access}) traced a connection-lifecycle fix applied to exactly one endpoint, \texttt{GET /genai/\allowbreak summary/\allowbreak\{article\_id\}}, after a real idle-in-transaction bug; every other handler in \texttt{articles.py}, \texttt{search.py}, and \texttt{genai.py}'s \texttt{chat} still closes its connection with no \texttt{try}/\texttt{finally} guard, and no handler anywhere pools connections rather than opening a fresh \texttt{psycopg2.connect(...)} per request. Introducing a shared connection pool (\texttt{psycopg2.pool} or an async driver) and applying the same \texttt{try}/\texttt{finally} discipline uniformly would remove both the leak risk and the per-request connection overhead at once.
    \item \textbf{Arabic-language content and Arabic-dialect sentiment.} Chapter 1 named that some of the same four outlets already publish in Arabic alongside French, but every entry in \texttt{RSS\_FEEDS} is tagged \texttt{"fr"}, and \texttt{cardiffnlp/twitter-xlm-roberta-base-sentiment} (Chapter 5, Section~\ref{sec:sentiment}), while multilingual, was neither selected nor validated for Tunisian Arabic or its \textit{Derja} dialect. Ingesting Arabic-language articles and either validating or fine-tuning a dialect-aware sentiment model would extend real coverage rather than only the multilingual embedding model's nominal capability.
    \item \textbf{Enabling the LLM region resolver.} Chapter 5 (Section~\ref{sec:enrichment}) and Chapter 2 both noted that \texttt{tag\_region\_llm} is fully implemented and tested in \texttt{etl/regions.py} but never exercised: \texttt{transform\_articles} calls \texttt{tag\_region} with \texttt{use\_llm=False}, so an article the gazetteer cannot place unambiguously is left \texttt{region = NULL} rather than falling through to the LLM tier. Turning that flag on, at least for articles the gazetteer left unresolved, would shrink the "National / non-localisé" bucket the map dashboard currently falls back to.
    \item \textbf{Named-entity extraction.} Chapter 4 (Section~\ref{sec:schema}) described the \texttt{entities} table -- one row per extracted named entity, cascading on its parent article -- as forward-provisioned schema with no pipeline step that ever inserts into it. Wiring a lightweight named-entity recognition step (a spaCy French pipeline, already a project dependency for BERTopic's stop-word list, or an LLM-based extractor following the same pattern as theme classification) into \texttt{transform\_articles} would put that table to its intended use.
    \item \textbf{Authentication for operator and pipeline endpoints.} Chapter 8 (Section~\ref{sec:backend-security}) was explicit that no endpoint in the API requires authentication or an API key, including \texttt{POST /pipeline/\allowbreak trigger} and \texttt{POST /pipeline/\allowbreak cluster} -- both reachable by anyone who can reach the API's port, currently workable only because that port is not the dashboard's real front door. Adding at least a shared operator token or API key on the pipeline and clustering routes, ahead of any future deployment where the API's port is more directly exposed, would close that gap before it becomes a real incident rather than a documented risk.
\end{itemize}
\end{sloppypar}

\section{Closing remarks}
\label{sec:closing-remarks}
\begin{sloppypar}
None of the eight gaps above was discovered by inspection after the fact; each was named, in the chapter where it belongs, at the moment the code that has the gap was itself being read. That is a deliberate property of this report and, more importantly, of the project it describes: a three-person, one-semester academic team building a Groq free-tier, single-database, five-service platform under real time and cost constraints, choosing a coarse-grained Prefect task over one task per pipeline step, a gazetteer before an LLM fallback, and a hardened connection-close path for one endpoint before generalizing it to five -- trade-offs made deliberately, and documented as such, rather than found later as unexplained shortcuts. News Insight, as it stands at the end of this report, turns four independent, French-language Tunisian news sources into one continuously updated, semantically searchable, sentiment- and topic-tagged corpus, explorable through a chat assistant that cites its sources and a governorate map built without a mapping library -- while carrying, openly, the list above as exactly what a next iteration should take on first.
\end{sloppypar}
